\section{Problem Formulation}
% In 1--2 paragraphs, describe the problem addressed in your deep learning project in a clear and concrete way. You may consider including the following:
% \begin{itemize}
% \item What task your project focuses on (e.g., image classification, object detection, semantic segmentation, speech recognition, time-series prediction, recommendation, or multi-modal learning).
% \item What information is available to the model, such as the input data, features, labels, context, or additional sources of information.
% \item What kind of supervision or training signal is used, such as labeled data, expert annotations, synthetic data, self-supervised learning, or multi-task objectives.
% \item What optimization objective your project aims to solve, including the loss function, evaluation target, and any practical constraints.
% \item What assumptions your project makes regarding the data, model design, or computational resources.
% \end{itemize}
This project focuses on an Image-to-Sequence translation task within the domain of Optical Music Recognition (OMR). We frame the transcription of standard staff notation into numbered musical notation (Jianpu) as an end-to-end sequence generation problem. The model takes 2D images of printed sheet music as input and directly decodes them into 1D arrays of custom semantic tokens representing Jianpu characters, rhythms, and performance marks. The training process utilizes fully supervised learning, guided by paired datasets of sheet music images and their corresponding ground-truth token sequences. To address the scarcity of such paired datasets, the training signal will primarily rely on large-scale synthetic data generated by rendering digital score files.

The primary optimization objective is to minimize the sequence-level cross-entropy loss between the predicted and ground-truth token sequences, with the ultimate goal of maximizing token-level accuracy and exact sequence match rates. To maintain project feasibility, we make a few key assumptions: first, the input data is strictly restricted to printed scores (e.g., clear scans or photographs of printed sheet music), explicitly excluding the highly variable domain of handwritten scores. Second, we assume that our custom tokenization vocabulary is sufficiently expressive to capture all essential musical semantics required for music performance without structural information loss.